\section{Mixed states, entanglement measures, and noise}
\hypertarget{mixed-states}{}

Every simulation in QPU passes through one \emph{physics engine}. The compiler
turns source into gates and the simulator decides the order in which they run,
but the physics engine alone changes the quantum state and computes the
quantities in this chapter. The particle cards in the Circuit builder show
several of them.

\subsection{Density matrices}

A pure state $\ket\psi$ can also be written as the \emph{density matrix}
\[
\rho=\ket\psi\bra\psi,\qquad \rho_{ij}=a_i\,\overline{a_j}.
\]
The diagonal entries are the basis probabilities. The off-diagonal entries,
called \emph{coherences}, carry the relative phases that make interference
possible. A \emph{mixed state} is a weighted mixture of pure states,
\[
\rho=\sum_k p_k\ket{\psi_k}\bra{\psi_k},\qquad p_k\ge0,\quad \sum_k p_k=1,
\]
and describes either classical uncertainty about which state was prepared or a
part of a larger entangled system. Every density matrix is Hermitian, has
trace 1, and has no negative eigenvalues. Global phase cancels in
$\ket\psi\bra\psi$, so a density matrix records exactly the physically
observable part of a state.

\subsection{Reduced states and the partial trace}

A single wire of an entangled register has no state vector of its own, but it
always has a density matrix. For a system $AB$ the state of $A$ alone is the
\emph{partial trace}
\[
\rho_A=\operatorname{Tr}_B(\rho_{AB}),\qquad
(\rho_A)_{ij}=\sum_{k}(\rho_{AB})_{(i,k),(j,k)}.
\]
It gives the correct statistics for every measurement made on $A$ alone. For
one qubit it fixes a \emph{Bloch vector} $(x,y,z)$ through
\[
\rho=\tfrac12\left(I+xX+yY+zZ\right),\qquad
x=\langle X\rangle,\quad y=\langle Y\rangle,\quad z=\langle Z\rangle.
\]
Pure states lie on the Bloch sphere ($r=1$); mixed states lie inside the
Bloch ball ($r<1$). The particle cards draw each wire from its reduced state,
which is why an entangled wire's arrow shrinks toward the centre.

\subsection{Purity and mixedness}

The purity of a state is
\[
P=\operatorname{Tr}(\rho^2)=\sum_{ij}|\rho_{ij}|^2 .
\]
$P=1$ exactly for pure states. For one qubit $P=(1+r^2)/2$, which ranges from
1 on the sphere to $\tfrac12$ at the centre (the maximally mixed state $I/2$).
The particle card's \emph{mixedness} is the normalized linear entropy
$M=2(1-P)$: 0 for a pure wire and 1 for a maximally mixed one.

\begin{beginnerbox}[Mixed does not mean noisy]
Without noise, the whole register is always in a pure state. A wire can still
look mixed because it is entangled with other wires: its information lives in
the correlations, not in the wire alone. When the register is pure and a wire
is mixed, the card shows \textbf{Entangled with register}.
\end{beginnerbox}

\subsection{Entropy and entanglement}

The \emph{von Neumann entropy} measures how mixed a state is, in bits:
\[
S(\rho)=-\operatorname{Tr}(\rho\log_2\rho)=-\sum_k\lambda_k\log_2\lambda_k,
\]
where $\lambda_k$ are the eigenvalues of $\rho$. A pure state has $S=0$; a
maximally mixed qubit has $S=1$.

If the whole register is pure, a subsystem $A$ is entangled with the rest
exactly when $\rho_A$ is mixed, and $S(\rho_A)$ is the \emph{entanglement
entropy}: how many bits of entanglement join $A$ to the rest.

\begin{center}
\begin{tabularx}{\textwidth}{l X X X}
\toprule
State & One qubit & Two qubits together & Whole register\\
\midrule
$\ket{00}$ & $P=1$, $S=0$ & $P=1$ & pure\\
Bell $(\ket{00}+\ket{11})/\sqrt2$ & $P=\tfrac12$, $S=1$ & $P=1$ & pure\\
GHZ $(\ket{000}+\ket{111})/\sqrt2$ & $P=\tfrac12$, $S=1$ & $P=\tfrac12$, $S=1$ & pure\\
\bottomrule
\end{tabularx}
\end{center}

The GHZ row shows why entanglement cannot be reduced to pairwise links. Any
two GHZ qubits are only classically correlated (00 or 11 with equal chance),
yet all three together form a single pure, entangled state. Questions about
entanglement are therefore always about a chosen split of the register into a
subsystem and the rest.

When the whole register is mixed, for example after noise, local mixedness no
longer proves entanglement. The physics engine then uses the \emph{partial
transpose} test: transpose only $B$'s indices of $\rho_{AB}$. If the result
has a negative eigenvalue the state is entangled, and the sum of the negative
eigenvalues' magnitudes is the \emph{negativity} ($\tfrac12$ for a Bell pair).
A negative eigenvalue is proof for any split. The converse holds only for two
qubits (or a qubit and a qutrit): for larger splits some entangled states have
no negative eigenvalue at all. The engine therefore reports one of three
results, and says which test it used:

\begin{center}
\begin{tabularx}{\textwidth}{l X}
\toprule
Result & Meaning\\
\midrule
entangled & Proven: the reduced state of a pure register is mixed, or the
partial transpose has a negative eigenvalue.\\
separable & Proven: the reduced state of a pure register is pure, or a
two-qubit mixed state has no negative eigenvalue.\\
inconclusive & Not decided: a larger mixed split passed the partial-transpose
test, or the register is too large to test. This is \emph{not} evidence that
the wires are unentangled.\\
\bottomrule
\end{tabularx}
\end{center}
A particle card shows \textbf{Entangled with register} only for a proven
result and \textbf{Entanglement undetermined} for an inconclusive one; it
never labels a wire ``not entangled'' because a test failed to find
entanglement.

QPU has no \code{ENTANGLE} instruction. Entanglement always comes from
ordinary gates such as \code{H} followed by \code{CNOT}; the physics engine
only detects and measures it.

\subsection{Measurement in other bases}

\code{MEASURE} measures in the computational (Z) basis. The physics engine can
also measure the X and Y observables by rotating the wire's eigenbasis onto Z,
measuring, and rotating back:
\begin{center}
\begin{tabular}{llll}
\toprule
Basis & Outcome 0 & Outcome 1 & Rotation onto Z\\
\midrule
Z & $\ket0$ & $\ket1$ & none\\
X & $\ket+$ & $\ket-$ & $H$\\
Y & $\ket{+i}$ & $\ket{-i}$ & $HS^\dagger$\\
\bottomrule
\end{tabular}
\end{center}
The state $\ket+$ gives a random result in the Z basis but always gives 0 in
the X basis: what a measurement reveals depends on the observable. In QPU
source, \code{MEASURE -I Q -BASIS X} (or \code{Y}) chooses the observable; a
\code{MEASURE} without \code{-BASIS} is a Z measurement.

\subsection{Fidelity}

\emph{Fidelity} measures how close a simulated state $\rho$ is to an intended
state $\sigma$:
\[
F(\rho,\sigma)=\left(\operatorname{Tr}\sqrt{\sqrt\rho\,\sigma\sqrt\rho}\right)^2,
\qquad
F(\ket\psi,\ket\phi)=|\langle\psi|\phi\rangle|^2 .
\]
$F=1$ means the states are identical (up to global phase) and $F=0$ means they
are perfectly distinguishable. For example, $F(\ket+,\ket0)=\tfrac12$.

\subsection{Reset as a physical operation}

RESET forces a wire to $\ket0$. As a density-matrix operation it is the
\emph{reset channel}: the wire ends in $\ket0$ and whatever it shared with the
other wires is traced out,
\[
\rho\;\longmapsto\;\ket0\bra0\otimes\operatorname{Tr}_q(\rho).
\]
A state vector cannot hold the mixed result of resetting an entangled wire,
so a state-vector simulation follows one run of this channel the way real
hardware does: measure the wire in Z, then apply X if the result was 1.
Averaged over many runs this gives exactly the reset channel. The
\hyperlink{reset-semantics}{Reset} chapter describes what this means for
circuits.

\subsection{Noise and decoherence}

Real qubits are not isolated. Their interaction with the environment is
described by \emph{quantum channels}
\[
\rho\;\longmapsto\;\sum_k K_k\,\rho\,K_k^\dagger,
\qquad \sum_k K_k^\dagger K_k=I,
\]
which turn pure states into mixed ones. The physics engine provides these
single-qubit channels:
\begin{center}
\begin{tabularx}{\textwidth}{l X}
\toprule
Channel & Effect\\
\midrule
Bit flip ($p$) & $\rho\mapsto(1-p)\rho+pX\rho X$: $\ket0$ and $\ket1$ swap with probability $p$.\\
Phase flip ($p$) & $\rho\mapsto(1-p)\rho+pZ\rho Z$: coherences shrink by $1-2p$.\\
Depolarizing ($p$) & $\rho\mapsto(1-p)\rho+p\,I/2$: the Bloch vector shrinks by $1-p$.\\
Amplitude damping ($\gamma$) & Energy loss: $\ket1$ decays to $\ket0$ with probability $\gamma$.\\
Phase damping ($\lambda$) & Coherences shrink by $\sqrt{1-\lambda}$; populations are unchanged.\\
\bottomrule
\end{tabularx}
\end{center}

\emph{Decoherence} is usually quoted as two times. $T_1$ is the relaxation
time: after time $t$ an excited qubit remains in $\ket1$ with probability
$e^{-t/T_1}$. $T_2$ is the coherence time: off-diagonal entries decay as
$e^{-t/T_2}$. Physics requires $T_2\le2T_1$.

\begin{cautionbox}[Cycles are not time]
QPU cycles (\code{INCREASECYCLE}) are logical program stages, not elapsed
physical time. Decoherence uses a separate timing model that assigns a
duration to each gate. Noise is always optional: ordinary runs are ideal and
use the cheaper state vector ($2^n$ amplitudes) rather than a density matrix
($4^n$ entries). Choose \textbf{Run with noise} in the \textbf{Physics
inspector} on the Particle visualization page to rerun a circuit with a noise
channel and optional $T_1$/$T_2$; the inspector reports the fidelity with the
ideal run. Noise settings are not part of QPU source.
\end{cautionbox}

\subsection{No cloning and teleportation}

CNOT copies a known basis value onto a $\ket0$ target, but applied to
$\ket+\ket0$ it produces a Bell pair, not $\ket+\ket+$. The two results have
fidelity $\tfrac12$ and the target becomes entangled rather than a copy. This
is the \emph{no-cloning theorem}: no operation can copy an unknown quantum
state.

Quantum state can still be \emph{moved}. SWAP exchanges two wires' states, and
teleportation transfers a state using only ordinary operations: share a Bell
pair (\code{H}, \code{CNOT}), entangle the message with one half (\code{CNOT},
\code{H}), measure two wires, and apply \code{X} and \code{Z} to the
destination under classical \code{-IF} conditions. The destination then holds
the original state with fidelity 1 and the source wire has been measured, so
the state was never in two places at once. QPU has no \code{CLONE} or
\code{TELEPORT} instruction.
